\section{Limitations and Conclusion}
\label{sec:conclusion}

\subsection{Limitations of this review}

This review is intentionally broad but not exhaustive. The archive was assembled
for a concrete research decision, which biases it toward dexterous hands,
learning-based control, recent tactile foundation models, and graph/contact
representations. It under-samples classical grasp mechanics, human haptics,
teleoperation, non-learning force control, and some industrial dual-arm systems.
Screening and coding were performed by one reviewer. Five records do not have a
locally validated PDF, and publication/code/data status can change rapidly.
Several 2026 methods are preprints or accepted/programme-listed works; their final
claims and artifacts may differ. Thirty purposively selected core papers received
full-text precision coding, while the remaining 43 scholarly records retain the
lighter archive-wide, rule-assisted coding layer. The precision subset is not a
random sample, and neither layer used dual independent extraction or inter-rater
agreement. Transparent evidence cards reduce ambiguity but do not replace those
controls or independent reproduction.

The corpus is too heterogeneous for meta-analysis. Reported success rates mix
different embodiments, sensors, training budgets, objects, horizons, and
definitions. Counts therefore describe only the curated archive. The review also
does not claim that load paths are the only useful abstraction. Deformation,
compliance, friction cones, contact fields, hybrid modes, and action-conditioned
latent dynamics may be superior in particular regimes.

\subsection{Conclusion}

Tactile robotics and bimanual manipulation have each advanced rapidly, but their
intersection exposes a missing representational layer. Local tactile encoders
identify what happens at a sensor; bimanual coordination requires a model of how
several interfaces become coupled through shared entities and how that coupling
changes at contact transitions. Graphs, spatial anchors, contact fields, modes,
and unstructured temporal attention are competing ways to express this layer.
None should be presumed best.

The strongest immediate benchmark is quasi-static handover because it provides
observable phases and controlled load redistribution. Cooperative tool--workpiece
manipulation is an established but complementary extension that adds a longer and
role-asymmetric load path. The narrower question of matched, observation-grounded
cross-hand load flow is comparatively sparse in the mapped corpus; it belongs
after the simpler mechanism has passed
matched-budget and coordination-necessity gates. Across both tasks, the decisive
methodological requirements are deployment-observable inputs, explicit label
provenance, episode/object/contact-sequence splits, oracle--inferred--random
structure comparisons, uncertainty, and interventions.

The full-text core reinforces this requirement: zero of 30 papers evaluated
strict held-out contact-topology transfer under deployment-observable inference.
This is a scoped evidence-map result, not an estimate of the entire literature,
and it motivates a decisive experiment rather than a priority claim.

An explicit bimanual contact-flow encoder can therefore be a strong contribution,
but the graph layer is not the contribution. The contribution is a falsifiable
factorization of sensed interfaces and object-mediated load propagation, together
with evidence that this factorization improves prediction, control, or transfer
under conditions where simpler temporal fusion does not. That framing provides a
clear starting point for experiments and a durable standard against which future
structured tactile representations can be judged.
